\section{Firmware Architecture and Implementation}
\label{sec:architecture}

This chapter walks through the firmware running on the EPS
microcontroller. It begins with the design philosophy that shapes
every decision below, then describes each major module --- the main
loop, the state machine, the MPPT algorithm, the PWM driver, the
sensor abstraction, the CHIPS codec, and the manual-control
mode --- and closes with the conventions that the code must comply
with.

\subsection{Design philosophy}

The firmware is bare-metal C running in a single super-loop. There
is no real-time operating system, no vendor hardware abstraction
layer, no dynamic memory allocation after boot, and no recursion.
Every loop or buffer has a compile-time-known upper bound. These
constraints come directly from the NASA Power of 10 rules
\cite{holzmann_powerof10_2006}, the JPL Institutional Coding
Standard for C \cite{jpl_d60411_2009}, and MISRA C:2004
\cite{misra_c_2004}, and they exist for the same reason any
flight-software project adopts them: behaviour that an auditor can
reason about statically is much easier to argue safe than behaviour
that depends on runtime memory pressure or scheduling order.

A second principle is the deliberate separation of \textit{pure
logic} from \textit{hardware access}. Functions that compute --- the
state machine dispatcher, the MPPT iteration, the CHIPS frame
codec, the CRC calculator --- never touch registers and never use
global state. They take input structs, return output structs, and
can be compiled on a laptop and tested in isolation. Functions
that touch hardware --- the SERCOM UART driver, the TCC PWM
driver, the ADC reader, the I\textsuperscript{2}C master --- never
contain control logic. This split makes every algorithmic claim of
this report testable without any chip on the bench.

\subsection{Folder and naming conventions}

The source tree is organised by responsibility, not by topic. The
top-level project folder \texttt{src-pds/} contains one subfolder
per architectural layer:

\begin{lstlisting}
src-pds/
  app/                              main loop
  board_command_contract/           CHIPS command IDs and payloads
  board_hardware_startup/           clock, peripherals, GPIO init
  board_outputs/                    PWM and output-state writes
  communication_with_esp32/         CHIPS read/dispatch/reply
  externally_controlled_behaviors/  per-mode runners
  runtime_state/                    shared snapshot RAM
  sensor_inputs/                    layered sensor abstraction
  shared_helpers/                   byte-order helpers
  state_machine_pure_logic/         pure dispatcher function
  status_reporting_to_esp32/        telemetry builders
\end{lstlisting}

File names are sentences, not abbreviations. A reader inspecting
the folder listing can tell what each file does without opening it:
\texttt{functions\_to\_run\_mppt\_algorithm\_with\_selected\_input\_source.c},
\texttt{functions\_to\_apply\_allowed\_pwm\_to\_board\_hardware.c},
\texttt{functions\_to\_initialize\_board\_hardware\_before\_main\_loop\_runs.c},
and so on. This style is taken directly from the project's
\texttt{conventions.md} and is enforced on every file.

\subsection{The main loop}
\label{sec:main-loop}

The application entry point reads like a table of contents. Listing
\ref{lst:main-loop} reproduces the body of \texttt{main()} verbatim
from \texttt{src-pds/app/main.c}, with only minor cosmetic
shortening.

\begin{lstlisting}[caption={The main super-loop of the EPS firmware.},label={lst:main-loop}]
int main(void) {
    switch_mcu_clock_from_1mhz_to_48mhz();
    setup_sercom0_uart_registers_for_esp32_link();
    setup_millisecond_timer_for_loop_timing();
    setup_sensor_reading_hardware();
    setup_pwm_output_starting_off();
    setup_output_pins_starting_off();

    start_esp32_command_reader_with_empty_message_state();
    set_starting_runtime_state_to_safe_defaults();

    while (1) /* @non-terminating@ */ {
        read_and_execute_commands_from_esp32();
        toggle_status_led_for_heartbeat_if_due_and_not_in_manual_mode();

        if      (requested_mode_is_off())            run_off_mode();
        else if (requested_mode_is_flight())         run_full_satellite_logic();
        else if (requested_mode_is_mppt_test())      run_mppt_test_only();
        else if (requested_mode_is_state_test())     run_state_transition_test_only();
        else if (requested_mode_is_fixed_pwm_test()) run_fixed_pwm_test_only();
        else if (requested_mode_is_manual())         run_manual_control_mode();

        block_dangerous_outputs();
        apply_outputs_to_board();
        send_status_if_needed();
        reset_watchdog_timer();
    }
}
\end{lstlisting}

The mode-dispatch block is one of the project's most explicit
design choices. Rather than building a single monolithic
``flight'' loop that runs every behaviour simultaneously, the
firmware exposes each capability as a separately commandable mode.
Each mode runs independently, in isolation from the others. This
keeps every individual capability testable; once each is verified,
the integration into a single ``flight'' loop becomes one
additional method, which is the
\texttt{run\_full\_satellite\_logic()} branch above.

\subsection{The state machine}

The PCU operating-mode state machine is implemented as a single
pure function in
\texttt{state\_machine\_pure\_logic/}\\
\texttt{functions\_to\_compute\_next\_pcu\_state\_and\_actuator\_commands\_from\_pure\_logic.c}.
The function signature takes a struct of sensor readings (battery
voltage, panel voltage, panel current, temperature, OBC heartbeat,
OBC-commanded mode, fault flags) and returns a struct of actuator
commands (selected PCU mode, safe-mode sub-state, MPPT enable
flag, heater enable flag, load mask, deploy-antenna flag,
send-telemetry flag, reboot flag).

The dispatcher reduces to a short cascade of \texttt{if}/\texttt{else
if} clauses: if any safety condition is tripped, return
\texttt{SAFE\_MODE}; otherwise return
\texttt{BATTERY\_DISCHARGE} (no sun), \texttt{MPPT\_CHARGE}
(below battery-full), \texttt{SA\_LOAD\_FOLLOW} (battery full),
or \texttt{CV\_FLOAT} (the intermediate band). The full flight FSM
includes time-based transitions (the \texttt{t1} and \texttt{t2}
debounce counters), hysteresis bands, and per-mode internal
algorithms; these are explicitly listed as future work in
Chapter~\ref{sec:results}.

A complementary verification function, the \texttt{dispatcher()},
is implemented separately and compared against the FSM output at
steady state. This is the standard
specification-versus-implementation cross-check that catches
regressions when the FSM is extended.

\subsection{The MPPT algorithm}

The MPPT module
(\texttt{src/mppt\_algorithm.c}) implements Incremental Conductance
in pure integer arithmetic. The algorithm reads the panel's raw
ADC values for voltage and current, filters them through an
8-sample moving average, computes the IncCond decision, and
returns a new duty cycle in the range $[0, 65535]$ that the PWM
module will apply to the buck converter.

\begin{lstlisting}[caption={Skeleton of the Incremental Conductance iteration.},label={lst:mppt}]
uint16_t mppt_algorithm_run_one_iteration(
        struct mppt_algorithm_state *state,
        uint16_t raw_panel_voltage_adc,
        uint16_t raw_panel_current_adc) {

    update_moving_average_filter(state, raw_panel_voltage_adc, raw_panel_current_adc);

    int32_t delta_v = state->filtered_v - state->previous_v;
    int32_t delta_i = state->filtered_i - state->previous_i;

    /* Decide whether to step duty up or down based on di/dv
     * and i/v comparison; clamp to [DUTY_MIN, DUTY_MAX]. */
    state->duty_cycle = apply_incremental_conductance_decision(
            state, delta_v, delta_i);

    state->previous_v = state->filtered_v;
    state->previous_i = state->filtered_i;

    return state->duty_cycle;
}
\end{lstlisting}

Because the function takes only the algorithm's own state struct
and the two ADC values, it can be run on a laptop with simulated
ADC inputs. A reference simulator (a five-parameter single-diode
panel model) was built for exactly this purpose during the dev-
board phase of the project; the same algorithm now runs unchanged
on the SAMD21.

\subsection{The PWM driver}

The PWM driver
(\texttt{src/drivers/pwm\_buck\_converter\_tcc0\_pa12\_pa13\_on\_mainboard.c})
configures TCC0 to produce two complementary outputs at 300\,kHz
on pins PA12 and PA13 with hardware-inserted dead-time. The
driver exposes one initialisation function and one runtime
function:

\begin{lstlisting}
void pwm_tcc0_initialize_complementary_300khz(void);
void pwm_tcc0_set_duty_cycle(uint16_t duty_as_fraction_of_65535);
\end{lstlisting}

The duty argument matches the MPPT output range, so the two modules
plug together without any rescaling. Internally, the driver maps
the 16-bit duty value to the TCC0 \texttt{CC} register range
(0--160 for 300\,kHz on the 48\,MHz GCLK), and writes both compare
channels under the dual-channel TCC0 dead-time scheme required by
the actual PA12/PA13 routing on the mainboard.

The driver also implements an explicit \textit{arm/disarm} safety
gate. PWM output is disarmed on boot, on every mode change, and on
any injected-fault flag. Re-arming requires an explicit
\texttt{pwm-arm} command from the host. This prevents the buck
stage from being unintentionally energised during firmware
development.

\subsection{The sensor abstraction}

The firmware reads sensor values through a three-layer
abstraction (Figure~\ref{fig:sensor-layers}). The state machine
and the MPPT algorithm at the top call exactly five functions
--- \texttt{read\_battery\_voltage()},
\texttt{read\_battery\_current()},
\texttt{read\_panel\_voltage()}, \texttt{read\_panel\_current()},
\texttt{read\_charging\_rail\_voltage()} --- and get back one
number each in engineering units. They never touch the ADC,
I\textsuperscript{2}C, or any specific chip. A single runtime
flag selects whether those numbers come from injected operator
values (the default during development) or from real PCB
hardware. The hardware path itself is split per chip
(INA226, LT6108, TPS25940 IMON, resistor dividers), each behind
its own small driver.

\begin{figure}[ht]
\centering
\fbox{\parbox{0.8\textwidth}{\centering [Figure placeholder]\\
Sensor abstraction in three layers: top layer = five
\texttt{read\_X()} functions called by state machine and MPPT;
middle layer = per-chip drivers (INA226, LT6108, TPS25940 IMON,
voltage dividers) plus the injected-value branch; bottom layer =
two bus drivers (ADC reader and I\textsuperscript{2}C master).}}
\caption{Three-layer sensor abstraction. The top layer hides the
sensor-source decision from the rest of the firmware.}
\label{fig:sensor-layers}
\end{figure}

This split is what lets the same firmware run on the dev board
(no real sensors --- everything falls back to the injected value)
and on the real mainboard (the same hardware path now drives
actual silicon).

\subsection{The CHIPS frame codec}

The CHIPS protocol is implemented as a pure-logic codec in
\texttt{chips\_protocol\_encode\_decode\_frames\_with\_crc16\_kermit.c}.
The encoder builds a framed byte stream from a (sequence, command,
payload) tuple; the decoder is a byte-fed state machine that
accumulates bytes between sync markers, runs the CRC, and emits
either \texttt{CHIPS\_FRAME\_READY}, \texttt{CHIPS\_INCOMPLETE},
or \texttt{CHIPS\_ERROR}. Byte stuffing is applied between the two
sync markers, exactly as specified by the mission document.

On the SAMD21 side, the codec is wrapped by a thin transport layer
that reads bytes from the SERCOM0 UART (interrupt-driven, with
256-byte ring buffers in both directions). On the ESP32 side, the
same codec runs in an Arduino sketch so that the test harness can
build valid frames from the laptop's plain-text commands.

\subsection{Manual control mode}

A separate operating mode lets an operator drive every user-facing
board output independently from a web page, without involving the
state machine or the MPPT loop. The mode is selected by the
\texttt{enter\_manual} CHIPS command; four subsequent commands
(\texttt{set\_manual\_pwm}, \texttt{set\_manual\_pv},
\texttt{set\_manual\_bat}, \texttt{set\_manual\_led}) write directly
to the corresponding actuator. This mode exists primarily for
hardware bring-up and one-off bench verifications; it is not part
of any flight scenario.

\subsection{Coding conventions}

Every file in the project is required to comply with the rules
documented in \texttt{conventions.md}. The most-binding rules
include: names are sentences not abbreviations; functions do one
job; assertions stay active in flight builds; zero compiler
warnings under \texttt{-Wall -Wextra -Werror -pedantic};
fixed-width integer types only; \texttt{volatile} on every
variable shared between an ISR and the main loop; one statement
per line; explicit parentheses in every compound expression. Each
of these rules has a stated rationale, and most are direct
translations of either the Power of 10 list
\cite{holzmann_powerof10_2006} or the JPL standard
\cite{jpl_d60411_2009}.

The total project compiles cleanly under
\texttt{-Wall -Wextra -Werror} for both the
\texttt{BOARD=devboard-pds} and \texttt{BOARD=mainboard-pds}
targets.
